\documentclass[letterpaper]{article}

\usepackage[margin=1.5cm]{geometry}
\usepackage{amsmath,amsfonts}
\usepackage{url}
\usepackage{bm}

\newcommand{\Real}{\mathbb{R}}
\newcommand{\tr}{^\top}
\newcommand{\opt}{^\star}

\setlength{\parskip}{3mm plus1mm minus1mm}
\setlength{\parindent}{0mm}

\title{Problem Set 9}
\author{Foundations of Machine Learning}
\date{}

\begin{document}
\maketitle

\begin{itemize}
\item \textbf{Submission}: Turn in as a \texttt{PDF} and a \textbf{separate source code file} on \url{http://mycourses.unh.edu}. Scanned hand-written notes are encouraged.
\item \textbf{Questions}: See the class website for a link to office hours
\end{itemize}

% Covers: Logistic regression (ISL 4, CO 7.1, MML 7.1) and Gradient descent for logistic regression (ISL 4, CO 9.3).

\textbf{Relevant reading}: ISL 4; CO 7.1, 9.3; MML 7.1. Problems marked \textbf{[code]} should be submitted as a program in a language of your choice (Julia or Python recommended). The cited exercise numbers are from ISL, 2nd edition (James et al., 2021).

\begin{enumerate}
\item ISL, Chapter 4, Conceptual exercises: problems 1, 6.
\item ISL, Chapter 4, Applied exercises: a logistic-regression exercise, e.g.\ on the \texttt{Weekly} data set \textbf{[code]}.
\item \textbf{Time spent.} Report approximately how many hours you spent completing this problem set (including the reading).
\end{enumerate}

\end{document}
